\phantomsection
\chapter*{Tragedi di Ruang Musik}
\addcontentsline{toc}{section}{Tragedi di Ruang Musik -- M. Ikhsan Fadilah}
\penulis{M. Ikhsan Fadillah}

\awalcerpen{A}{C di lantai tiga gedung} SMA Angkasa itu emang dari minggu lalu udah meledak mati-total, bikin aroma keringat nyampur tanah basah habis hujan sore kerasa banget di dalam ruang musik. Yulia asal selonjoran di karpet tipis dekat meja, mukanya ditutupi lembaran lirik yang udah lecek kena tumpahan es teh.

``Yul, awas itu kabel \textit{bass}-ku ketindih,'' tegur Maya sambil ngutak-ngatik kenop \textit{sound amplifier}. Suara kretek-kretek kasar keluar dari \textit{speaker}.

``Bentar, May$\dots$ mataku berat banget. Gara-gara Pak Budi ngasih soal matematika tadi buset dah, tiga nomor aja rasanya bikin otak mau pecah,'' keluh Yulia dari balik kertas leceknya.

Rina yang lagi ngetes \textit{double pedal drum}-nya mendadak hantam simbal kencang banget. ``Cshhh!''

Yulia langsung terlonjak bangun, kertas di mukanya kabur entah ke mana.

``Latihan, woi! Nggak ada mager-mageran!'' Rina ngomong sambil melotot, stik \textit{drum} di tangan kanannya ditunjuk-tunjuk ke arah Yulia sama Maya. ``Gila kalian ya, minggu depan itu pensi terakhir kita bertiga sebelum lulus. Aku nggak mau ya di panggung malah fals atau ada yang lupa \textit{part} lagu!!!''

Pintu ruangan berdecit pelan. Annisa, anak kelas sepuluh yang paling junior di antara mereka, nyelonong masuk sambil nggendong tas gitar yang tingginya hampir sama kayak badannya. ``Permisi, Kak$\dots$ maap telat, tadi piket kelas dulu.''

Tania yang dari tadi anteng di sudut ruangan cuma geleng-geleng kepala. Dia naruh lima cangkir plastik isi teh hangat sama sekantong plastik gorengan yang masih ngebul di meja. ``Udah, minum dulu. Rina juga jangan teriak-teriak terus, ntar suaramu habis duluan sebelum nyanyi.''

Mata Rina langsung tertuju ke bakwan jagung di meja. Sikap galaknya hilang seketika.

Mereka akhirnya ngumpul melingkar di lantai kayu yang agak berdebu. Obrolan ngalor-ngidul santai mengalir gitu aja. Dari ngebahas gosip anak IPA sampai lelucon konyol soal Pak Budi yang sering lupa naruh kacamata padahal ditaruh di jidatnya sendiri. Buat anak-anak ini, ruang musik berukuran lima kali empat meter ini bukan cuma tempat buat genjrengan, tapi tempat pelarian paling aman dari mumetnya tugas sekolah.

Jam dinding menunjukkan angka lima sore lewat sedikit. Langit di luar jendela kaca berubah jadi agak keabu-abuan. Latihan resmi akhirnya dimulai.

Lagu pertama tempo cepat meluncur mulus. Pukulan \textit{drum} Rina mantap, dentuman \textit{bass} Maya ngisi ritem dengan rapi, dan jemari Tania di \textit{keyboard} bikin suasana makin bernyawa. Yulia berdiri di tengah ruangan, genjrengan gitar merahnya kedengaran pas banget nyatu sama gitar Annisa.

Pas masuk lagu kedua yang temponya lebih pelan tapi emosional, suasana ruangan mendadak serius.

Masuk ke bagian \textit{interlude}, Yulia maju satu langkah mendekati \textit{amplifier}. Dia merem, benar-benar ngehayatin bagian \textit{solo} gitar yang dia ulik semalaman di rumah. Tangannya menekan kuat \textit{fret} belasan, mau ngambil teknik ``\textit{bend}'' senar nomor satu sampai nada tertinggi.

PLEKKK!!, Suara cetekan kawat putus berbunyi kencang banget, disusul bunyi cambukan keras di udara.

Senar gitar nomor satu yang tegangnya maksimal itu putus persis di dekat \textit{bridge}. Ujung kawat bajanya membal cepat ke atas, menghantam wajah Yulia.

``Aakh!'' Yulia menjerit kencang, refleks ngelepas gitar merahnya sampai jatuh ngegubrak ke lantai kayu.

Suara instrumen langsung mati total. Ruangan hening seketika.

``Yul! Kenapa Yul?!!!'' Maya langsung melompat ngelepas \textit{strap bass}-nya dan lari menubruk Yulia. Rina hampir menabrak \textit{cymbal}-nya sendiri pas lari muterin \textit{set drum}.

Yulia nunduk dalam-dalam sambil menutup mata kirinya pakai kedua telapak tangan. Darah segar mengalir cepat dari sela-sela jarinya, menetes-netes deras ke lantai kayu dan mengotori kerah seragam putihnya.

``Aduhh$\dots$ sakit banget, May$\dots$ sakit$\dots$'' erang Yulia. Suaranya gemetaran parah, napasnya patah-patah.

Maya histeris pas ngeliat merah darah udah memenuhi tangan Yulia. ``Jangan diusap! Yul, tanganmu jangan ditekan ke mata! Astaga, Rin, darahnya banyak banget!''

Tania panik setengah mati sampai cangkir plastik di meja tersenggol tumpah. ``Annisa!!! P3K di dekat pintu! Rina, cari Pak Satpam sekarang! Minta tolong dipanggilkan mobil atau apa aja!''

Rina tanpa pikir panjang langsung buka pintu lebar-lebar dan lari terbirit-birit di koridor lantai tiga yang udah mulai gelap sambil berteriak minta tolong. Annisa yang bongkar-bongkar kotak P3K tangannya gemetar hebat sampai isi kasa dan alkohol berantakan di lantai.

Di tengah ruangan yang mendadak terasa dingin dan mencekam itu, Yulia cuma bisa nangis sambil nahan nyeri luar biasa di mukanya. Maya memeluk bahu Yulia erat-erat, tidak peduli baju seragamnya sendiri mulai ikut kena noda darah. Sore yang harusnya penuh tawa itu hancur total dalam hitungan detik.

\jedahias
